% Part: first-order-logic
% Chapter: completeness
% Section: identity

\documentclass[../../../include/open-logic-section]{subfiles}

\begin{document}

\olfileid{fol}{com}{ide}
\olsection{તાદાત્મ્ય}

\begin{explain}
પાછલા વિભાગમાં આપેલી પદ-નિદર્શની રચના~$\eq$ ન ધરાવતા ગણો~$\Gamma$
માટે પ્રથમ-ક્રમ તર્કશાસ્ત્રની પૂર્ણતા સ્થાપવા પૂરતી છે. પદ-નિદર્શ
~$\eq$ ન ધરાવતા દરેક $!A \in \Gamma^*$ને સંતોષે છે (અને તેથી બધા
$!A \in \Gamma$ને પણ). પરંતુ~$\eq$ હાજર હોય ત્યારે આ રીત કામ
કરતી નથી. તેનું કારણ એ છે કે પછી~$\Gamma^*$માં કોઈ
!!a{sentence}~$\eq[t][t']$ હોઈ શકે, જ્યારે પદ-નિદર્શમાં કોઈ પણ પદનું
મૂલ્ય એ પદ પોતે જ છે. તેથી જો~$t$ અને~$t'$ જુદાં પદો હોય, તો
પદ-નિદર્શમાં તેમનાં મૂલ્યો—અનુક્રમે~$t$ અને~$t'$—પણ જુદાં છે અને આમ
$\eq[t][t']$ ખોટું છે. જોકે ``ભાગફલન'' તરીકે ઓળખાતી રચના વાપરીને આ
સમસ્યા ઉકેલી શકાય છે.
\end{explain}

\begin{defn}
  ધારો કે~$\Gamma^*$ એ~$\Lang L$નાં વાક્યોનો સુસંગત અને !!{complete}
  ગણ છે. ~ $\Lang L$નાં બંધ પદોના ગણ પર સંબંધ~$\approx$ આ રીતે
  વ્યાખ્યાયિત કરીએ છીએ:
  \[
  t \approx t' \text{\quad જો અને માત્ર જો \quad} \eq[t][t'] \in \Gamma^*
  \]
\end{defn}

\begin{prop}
\ollabel{prop:approx-equiv}
સંબંધ~$\approx$ને નીચેના ગુણધર્મો છે:
\begin{enumerate}
\item $\approx$ સ્વવાચક છે.
\item $\approx$ સમમિત છે.
\item $\approx$ પરંપરિત છે.
\item જો $t\approx t'$, $f$ !!a{function} અને $t_1$, \dots,
  $t_{i-1}$, $t_{i+1}$, \dots, $t_n$ બંધ પદો હોય, તો
\[
\Atom{f}{t_1,\dots,t_{i-1},t,t_{i+1},\dots,t_n} \approx
\Atom{f}{t_1,\dots,t_{i-1},t',t_{i+1},\dots,t_n}.
\]
\item જો $t\approx t'$, $R$ !!a{predicate} અને $t_1$, \dots,
  $t_{i-1}$, $t_{i+1}$, \dots, $t_n$ બંધ પદો હોય, તો
\begin{multline*}
\Atom{R}{t_1,\dots,t_{i-1},t,t_{i+1},\dots,t_n}\in\Gamma^*
\text{ જો અને માત્ર જો } \\
\Atom{R}{t_1,\dots,t_{i-1},t',t_{i+1},\dots,t_n}\in\Gamma^*.
\end{multline*}
\end{enumerate}
\end{prop}

\begin{proof}
~$\Gamma^*$ સુસંગત અને !!{complete} હોવાથી $\eq[t][t']\in\Gamma^*$ જો
અને માત્ર જો $\Gamma^*\Proves\eq[t][t']$. તેથી નીચેનું બતાવવું પૂરતું
છે:
\begin{enumerate}
\item દરેક બંધ પદ~$t$ માટે $\Gamma^*\Proves\eq[t][t]$.
\item જો $\Gamma^*\Proves\eq[t][t']$, તો $\Gamma^*\Proves\eq[t'][t]$.
\item જો $\Gamma^*\Proves\eq[t][t']$ અને $\Gamma^*\Proves
  \eq[t'][t'']$, તો $\Gamma^*\Proves\eq[t][t'']$.
\item જો $\Gamma^*\Proves\eq[t][t']$, તો
\[
\Gamma^* \Proves
\eq[\Atom{f}{t_1,\dots,t_{i-1},t,t_{i+1},\dots,t_n}][\Atom{f}{t_1,\dots,t_{i-1},t',t_{i+1},\dots,t_n}]
\]
દરેક $n$-સ્થાનીય !!{function}~$f$ અને બંધ પદો $t_1$, \dots,
$t_{i-1}$, $t_{i+1}$, \dots,~$t_n$ માટે.
\item જો $\Gamma^*\Proves\eq[t][t']$ અને
$\Gamma^*\Proves\Atom{R}{t_1,\dots,t_{i-1},t,t_{i+1},\dots,t_n}$,
તો દરેક $n$-સ્થાનીય !!{predicate}~$R$ અને બંધ પદો $t_1$, \dots,
$t_{i-1}$, $t_{i+1}$, \dots,~$t_n$ માટે
$\Gamma^*\Proves\Atom{R}{t_1,\dots,t_{i-1},t',t_{i+1},\dots,t_n}$.
\end{enumerate}
\end{proof}

\sourcecorrection{OLCO-011}{મૂળની \texttt{identity.tex}, પંક્તિ~69માં
પહેલા સંયુક્ત પદની દલીલ-યાદીમાં $t_{i+1}$ પછી વધારાનો અલ્પવિરામ છે.
અહીં તે દૂર કર્યો છે.}

\begin{prob}
\olref[fol][com][ide]{prop:approx-equiv}ની સાબિતી પૂરી કરો.
\end{prob}

\begin{defn}
ધારો કે~$\Gamma^*$ ભાષા~$\Lang L$નો સુસંગત અને !!{complete} ગણ છે,
~$t$ બંધ પદ છે અને~$\approx$ પાછલી વ્યાખ્યામાંનો સંબંધ છે. ત્યારે:
\[
\equivrep{t}{\approx}=\Setabs{t'}{t'\in\Trm[L],t\approx t'}
\]
અને
$\equivclass{\Trm[L]}{\approx}=\Setabs{\equivrep{t}{\approx}}{t\in\Trm[L]}$.
\end{defn}

\begin{defn}
\ollabel{defn:term-model-factor}
ધારો કે~$\Struct M=\Struct M(\Gamma^*)$ એ \olref[mod]{defn:termmodel}નો
~$\Gamma^*$ માટેનો પદ-નિદર્શ છે. ત્યારે
~$\Struct{\equivclass{M}{\approx}}$ નીચેની !!{structure} છે:
\begin{enumerate}
\item $\Domain{\equivclass{M}{\approx}}=\equivclass{\Trm[L]}{\approx}$.
\item $\Assign{c}{\equivclass{M}{\approx}}=\equivrep{c}{\approx}$.
\item $\Assign{f}{\equivclass{M}{\approx}}(\equivrep{t_1}{\approx},\dots,
  \equivrep{t_n}{\approx})=\equivrep{\Atom{f}{t_1,\dots,t_n}}{\approx}$.
\item $\tuple{\equivrep{t_1}{\approx},\dots,\equivrep{t_n}{\approx}}
  \in\Assign{R}{\equivclass{M}{\approx}}$ જો અને માત્ર જો
  $\Sat{M}{\Atom{R}{t_1,\dots,t_n}}$, એટલે કે, જો અને માત્ર જો
  $\Atom{R}{t_1,\dots,t_n}\in\Gamma^*$.
\end{enumerate}
\end{defn}

\begin{explain}
નોંધો કે આપણે~$\equivclass{\Trm[L]}{\approx}$નાં ઘટકો માટે
$\Assign{f}{\equivclass{M}{\approx}}$ અને
$\Assign{R}{\equivclass{M}{\approx}}$ને તેમનો ઉલ્લેખ
$\equivrep{t}{\approx}$ તરીકે કરીને, એટલે કે \emph{પ્રતિનિધિઓ}
~$t\in\equivrep{t}{\approx}$ દ્વારા વ્યાખ્યાયિત કર્યાં છે. આ વ્યાખ્યાઓ
પ્રતિનિધિઓની પસંદગી પર આધાર રાખતી નથી તેની ખાતરી કરવી પડશે: એટલે કે,
એ જ સામ્ય વર્ગો નક્કી કરતી બીજી કોઈ પસંદગીઓ~$t'$ માટે
($\equivrep{t}{\approx}=\equivrep{t'}{\approx}$) વ્યાખ્યાઓ એ જ પરિણામ
આપે છે. દાખલા તરીકે, જો~$R$ એક-સ્થાનીય !!{predicate} હોય, તો
વ્યાખ્યાની છેલ્લી કલમ કહે છે કે $\equivrep{t}{\approx}\in
\Assign{R}{\equivclass{M}{\approx}}$ જો અને માત્ર જો
$\Sat{M}{\Atom{R}{t}}$. જો $t\approx t'$ હોય એવા બીજા કોઈ પદ~$t'$ માટે
$\Sat/{M}{\Atom{R}{t'}}$, તો વ્યાખ્યા મુજબ
$\equivrep{t'}{\approx}\notin\Assign{R}{\equivclass{M}{\approx}}$ હોવું
પડે. જો $t\approx t'$, તો $\equivrep{t}{\approx}=
\equivrep{t'}{\approx}$; પરંતુ એક જ સમયે $\equivrep{t}{\approx}\in
\Assign{R}{\equivclass{M}{\approx}}$ અને $\equivrep{t}{\approx}\notin
\Assign{R}{\equivclass{M}{\approx}}$ હોઈ શકે નહિ. જોકે
\olref{prop:approx-equiv} ખાતરી આપે છે કે આવું બની શકતું નથી.
\end{explain}

\sourcecorrection{OLCO-012}{મૂળની \texttt{identity.tex}, પંક્તિ~134માં
બીજા પ્રતિનિધિ~$t'$ની ચર્ચા છતાં અસંતોષનું સૂત્ર
$\Sat/{M}{\Atom{R}{t}}$ લખાયું છે. દલીલને જોઈતું અને પછીના
નિષ્કર્ષમાં વપરાતું સૂત્ર $\Sat/{M}{\Atom{R}{t'}}$ હોવાથી અહીં એ
પુનઃસ્થાપિત કર્યું છે.}

\begin{prop}
~$\Struct{\equivclass{M}{\approx}}$ સુવ્યાખ્યાયિત છે; એટલે કે, જો
$t_1$, \dots, $t_n$, $t_1'$, \dots, $t_n'$ બંધ પદો હોય અને દરેક સૂચક
માટે $t_i\approx t_i'$ હોય, તો
\begin{enumerate}
\item $\equivrep{\Atom{f}{t_1,\dots,t_n}}{\approx}=
    \equivrep{\Atom{f}{t_1',\dots,t_n'}}{\approx}$, એટલે કે,
  \[
  \Atom{f}{t_1,\dots,t_n}\approx\Atom{f}{t_1',\dots,t_n'}
  \]
  અને
\item $\Sat{M}{\Atom{R}{t_1,\dots,t_n}}$ જો અને માત્ર જો
  $\Sat{M}{\Atom{R}{t_1',\dots,t_n'}}$, એટલે કે,
  \[
    \Atom{R}{t_1,\dots,t_n}\in\Gamma^* \text{ જો અને માત્ર જો }
    \Atom{R}{t_1',\dots,t_n'}\in\Gamma^*.
  \]
\end{enumerate}
\end{prop}

\begin{proof}
~$n$ પરના અનુમાન વડે \olref{prop:approx-equiv}માંથી ફલિત થાય છે.
\end{proof}

પદ-નિદર્શના કિસ્સાની જેમ, સત્યતા સહાયક પ્રમેય સાબિત કરતાં પહેલાં
આગળનું સહાયક પ્રમેય જરૂરી છે.

\begin{lem}
\ollabel{lem:val-in-termmodel-factored} ધારો કે
$\Struct M=\Struct M(\Gamma^*)$. તો
$\Value{t}{\equivclass{M}{\approx}}=\equivrep{t}{\approx}$.
\end{lem}
\begin{proof}
સાબિતી \olref[mod]{lem:val-in-termmodel}ની સાબિતી જેવી જ છે.
\end{proof}

\begin{prob}
\olref[fol][com][ide]{lem:val-in-termmodel-factored}ની સાબિતી પૂરી કરો.
\end{prob}

\begin{lem}
\ollabel{lem:truth}
દરેક વાક્ય~$!A$ માટે $\Sat{\equivclass{M}{\approx}}{!A}$ જો અને માત્ર
જો $!A\in\Gamma^*$.
\end{lem}

\begin{proof}
~$!A$ પર અનુમાન વડે, \olref[mod]{lem:truth}ની સાબિતીમાં હતું તેમ જ.
ફક્ત $!A\ident\eq[t][t']$ હોય તે કિસ્સા પર વધુ ધ્યાન આપવું પડે છે.
\begin{align*}
\Sat{\equivclass{M}{\approx}}{\eq[t][t']} & \text{ જો અને માત્ર જો }
\equivrep{t}{\approx}=\equivrep{t'}{\approx}
\text{ (\olref{lem:val-in-termmodel-factored} અને તાદાત્મ્યની વ્યાખ્યા મુજબ)}\\
& \text{ જો અને માત્ર જો }t\approx t'
\text{ ($\equivrep{t}{\approx}$ની વ્યાખ્યા મુજબ)}\\
& \text{ જો અને માત્ર જો }\eq[t][t']\in\Gamma^*
\text{ ($\approx$ની વ્યાખ્યા મુજબ).}
\end{align*}
\end{proof}

\sourcecorrection{OLCO-013}{મૂળની \texttt{identity.tex}, પંક્તિ~188માં
પહેલી સમતુલ્યતાને માત્ર ભાગફલ સંરચનાની વ્યાખ્યાથી મળતી કહે છે. પદોના
મૂલ્યો તેમના સામ્ય વર્ગો છે એ માટે તરત પહેલાંનું
\olref{lem:val-in-termmodel-factored} પણ જરૂરી છે; અહીં એ સંદર્ભ સ્પષ્ટ
કર્યો છે.}

\begin{digress}
નોંધો કે~$\Struct{M(\Gamma^*)}$ હંમેશાં !!{enumerable} અને અનંત છે,
પરંતુ~$\Struct{\equivclass{M}{\approx}}$ સાન્ત હોઈ શકે છે, કારણ કે
સામ્ય વર્ગો~$\equivrep{t}{\approx}$ સાન્ત સંખ્યામાં જ હોવાનું બની શકે
છે. આ અપેક્ષિત છે, કારણ કે~$\Gamma$માં એવાં !!{sentence}s હોઈ શકે જે
તેમને સાચાં બનાવતી કોઈ પણ !!{structure} સાન્ત હોવી જરૂરી બનાવે. દાખલા
તરીકે, $\lforall[x][\lforall[y][x=y]]$ સુસંગત !!{sentence} છે, પણ તે
ફક્ત બરાબર એક !!{element} ધરાવતા !!a{domain}વાળી !!{structure}sમાં જ
સંતોષાય છે.
\end{digress}

\end{document}
